\section{Finite-complement separation}\label{sec:finite-separation}

For a normal term $x\in F(D)$, let $\operatorname{Sub}(x)$ denote its finite
set of normal subterms.  If $X\subseteq F(D)$ is finite, set
\[
 S_X=\bigcup_{x\in X}\operatorname{Sub}(x),
 \qquad
 N_X=S_X\setminus A.
\]
Thus $N_X$ is finite and $S_X$ is closed under taking subterms.

\begin{lemma}[Finite-complement Rees quotient]
\label{lem:finite-complement-rees}
For every finite $X\subseteq F(D)$ there is a congruence $\rho_X$ on $F(D)$
such that every element of $A\cup N_X$ is a singleton $\rho_X$-class and, if
nonempty, $F(D)\setminus(A\cup N_X)$ is one additional class.  Consequently
$F(D)/\rho_X$ is a finite-complement compatible extension of $D$ and the
quotient map is injective on $X$.
\end{lemma}

\begin{proof}
Put
\[
 I_X=F(D)\setminus(A\cup N_X).
\]
We first observe that $I_X$ is absorbing in every coordinate: if one argument
of a basic operation lies in $I_X$, then so does the result.  Indeed, let
$f(s,t)$ be the result of applying a basic operation and suppose, for example,
that $s\in I_X$.  The result cannot lie in $A$, because reduction to a base
value is possible only when both arguments lie in $A$.  Nor can it lie in
$N_X$: if the normal term $f(s,t)$ were a selected subterm, then its subterm
$s$ would belong to $S_X$, contrary to $s\in I_X$.  The same argument applies
when $t\in I_X$.

Define
\[
 \rho_X=\Delta_{F(D)}\cup(I_X\times I_X).
\]
The absorption property makes $\rho_X$ a congruence; when $I_X$ is nonempty,
this is the Rees congruence determined by the absorbing subalgebra $I_X$
\cite{Tichy1981}.  Every base element remains a singleton class.  Outside the
image of $A$ the quotient has the singleton classes indexed by $N_X$, together
with at most one additional class represented by $I_X$.  Hence
\[
 |(F(D)/\rho_X)\setminus A|\le |N_X|+1<\infty.
\]
Since $X\subseteq A\cup N_X$, the quotient map is injective on $X$.
\end{proof}

The construction is the algebraic form of the familiar finite-subterm
construction with one sink state in tree automata \cite{ComonEtAl2008}.

\begin{corollary}[Separation]
\label{cor:finite-complement-separation}
If $x\ne y$ in $F(D)$, then there is a finite-complement compatible extension
$T$ for which the canonical map $F(D)\to T$ separates $x$ and $y$.
\end{corollary}

\begin{remark}
When $A$ is infinite, finite-complement extensions are not in general closed
under ordinary products over the diagonal copy of $A$: already the diagonal
copy of $A$ in $(A\sqcup\{\ast\})^2$ has infinitely many points outside it.
This observation is not needed in the proof above; it merely explains why the
usual product-of-separators argument does not stay inside the same class.
\end{remark}
